% !TEX root = ../Monografia.tex
% ---------------------------------------------------------------------------------------------------------------------------------------------------------
\section{Etapas de Desenvolvimento}
% ---------------------------------------------------------------------------------------------------------------------------------------------------------

O desenvolvimento do pacote \texttt{collector\_flutter} foi conduzido de forma iterativa e incremental, seguindo princípios de engenharia de software experimental conforme orientam Wohlin \textit{et al.}~\cite{wohlin2012} e Pressman e Maxim~\cite{pressman2020engenharia}.
As etapas foram planejadas de modo a integrar o rigor técnico com a experimentação empírica, garantindo a rastreabilidade entre objetivos, decisões de projeto e resultados obtidos.
O cronograma detalhado do TCC 1 e do TCC 2 foi preservado no Apêndice~\ref{ap:cronogramas}, pois funciona como registro de planejamento e execução do trabalho sem interromper a leitura dos procedimentos metodológicos centrais.

\subsection{Fase 1 – Revisão Bibliográfica e Fundamentação Teórica}

Na primeira fase, foi realizada uma revisão bibliográfica estruturada sobre métricas de desempenho, consumo energético e práticas de otimização em aplicações móveis.
A revisão foi organizada por eixos temáticos, de modo a identificar lacunas na literatura e fundamentar decisões de projeto com base em evidências verificáveis.

Foram consultadas bases de dados como IEEE Xplore, ACM Digital Library e SpringerLink, utilizando termos relacionados a \textit{software profiling}, \textit{mobile performance}, \textit{Flutter metrics}, e \textit{energy efficiency}.
Essa etapa permitiu consolidar o arcabouço teórico necessário à definição dos indicadores de desempenho a serem monitorados e das heurísticas de recomendação incorporadas ao pacote.
